\section{Fórmula de Tweedie}

\subsection{Enunciado de la fórmula e intuición}

La fórmula de Tweedie es una herramienta matemática clave para comprender la conexión entre la función score y los modelos de difusión. Responde a la siguiente pregunta: si conocemos la distribución marginal $p(x)$ y la distribución de verosimilitud $p(x \mid \mu)$, ¿podemos inferir la esperanza del variable latente $\mu$ a partir de la observación $x$?

\begin{importantbox}{Fórmula de Tweedie}
Considere el siguiente proceso generativo: primero se muestrea un valor medio $\mu$ desde una distribución a priori $p(\mu)$, y luego se muestrea una observación $x$ desde la distribución de verosimilitud $p(x \mid \mu) = \mathcal{N}(x;\; \mu, \sigma^2 I)$. Sea la distribución marginal $p(x) = \int p(x \mid \mu)\, p(\mu)\, d\mu$; entonces, el valor medio posterior es:
$$\mathbb{E}[\mu \mid x] = x + \sigma^2 \nabla_x \log p(x)$$
\end{importantbox}

\begin{knowledgebox}{Intuición de la fórmula de Tweedie}
Esta fórmula posee una explicación intuitiva muy elegante:
\begin{itemize}
    \item Si solo conocemos la observación $x$ y no tenemos otra información, la estimación más ingenua de $\mu$ es $x$ mismo.
    \item Pero la función score $\nabla_x \log p(x)$ proporciona una ``dirección de corrección'': apunta hacia regiones de alta densidad de probabilidad.
    \item La fórmula de Tweedie nos indica que la estimación óptima del valor medio posterior consiste en realizar un paso de corrección a partir de la estimación ingenua $x$, siguiendo la dirección del score, con una amplitud de corrección proporcional a la varianza del ruido $\sigma^2$.
\end{itemize}
\end{knowledgebox}

\subsection{Verificación en el caso gaussiano simple}

El ponente utiliza un ejemplo sencillo para verificar la validez y la intuición de la fórmula de Tweedie:

Sea $p(\mu) = \mathcal{N}(0, 1)$ y $p(x \mid \mu) = \mathcal{N}(x;\; \mu, \sigma^2)$. Dado que tanto la distribución a priori como la de verosimilitud son gaussianas, la distribución marginal también lo será (propiedad de las distribuciones gaussianas), por lo que:

$$p(x) = \mathcal{N}(0, 1 + \sigma^2)$$

Por tanto, la función score es:

$$\nabla_x \log p(x) = -\frac{x}{1 + \sigma^2}$$

Sustituyendo en la fórmula de Tweedie:

$$\mathbb{E}[\mu \mid x] = x + \sigma^2 \cdot \left(-\frac{x}{1+\sigma^2}\right) = x \cdot \frac{1}{1+\sigma^2}$$

\begin{knowledgebox}{Verificación intuitiva: estimación con contracción}
El resultado $\mathbb{E}[\mu \mid x] = \frac{x}{1+\sigma^2}$ presenta un efecto de ``contracción'' (shrinkage):
\begin{itemize}
    \item Cuando $x > 0$, $\mathbb{E}[\mu \mid x] < x$: la estimación del valor medio está más cerca de cero que el valor observado (en la dirección del valor medio a priori).
    \item Cuando $x < 0$, $\mathbb{E}[\mu \mid x] > x$: análogamente, se contrae hacia cero.
    \item Cuanto mayor es $\sigma^2$ (mayor es el ruido), más fuerte es la contracción: en presencia de alto ruido, confiamos más en la distribución a priori.
    \item Cuando $\sigma^2 \to 0$, $\mathbb{E}[\mu \mid x] \to x$: cuando desaparece el ruido, la observación misma constituye la mejor estimación.
\end{itemize}
Esto concuerda perfectamente con la intuición de la inferencia bayesiana.
\end{knowledgebox}

\subsection{Aplicación en modelos de difusión}

Al aplicar la fórmula de Tweedie a la distribución del salto hacia adelante en los modelos de difusión, $q(x_t \mid x_0) = \mathcal{N}(\sqrt{\bar{\alpha}_t}\, x_0, (1-\bar{\alpha}_t)I)$, donde el papel de ``valor medio'' lo desempeña $\sqrt{\bar{\alpha}_t}\, x_0$ y la ``varianza'' es $\sigma_t^2 = 1-\bar{\alpha}_t$, obtenemos:

$$\mathbb{E}[\sqrt{\bar{\alpha}_t}\, x_0 \mid x_t] = x_t + (1-\bar{\alpha}_t)\, \nabla_{x_t} \log q(x_t)$$

Reordenando los términos:

\begin{importantbox}{Forma de la fórmula de Tweedie en modelos de difusión}
$$\mathbb{E}[x_0 \mid x_t] = \frac{1}{\sqrt{\bar{\alpha}_t}} \left(x_t + (1-\bar{\alpha}_t)\, \nabla_{x_t} \log q(x_t)\right)$$
Esto indica que conocer la función score marginal $\nabla_{x_t} \log q(x_t)$ es equivalente a conocer la estimación óptima de $x_0$ dado $x_t$. La función score contiene toda la información necesaria para el ``desenfoque'' (denoising).
\end{importantbox}

\begin{warningbox}{Diferencia entre score condicional y score marginal}
El $\nabla_{x_t} \log q(x_t)$ que aparece aquí es el score de la distribución \textbf{marginal}, no el score de la distribución condicional $q(x_t \mid x_0)$. El score condicional se puede calcular directamente utilizando la fórmula de la distribución gaussiana, pero el score marginal requiere promediar mediante integración sobre todos los posibles $x_0$. Durante el entrenamiento, las redes neuronales aprenden implícitamente el score marginal mediante el promedio de grandes cantidades de muestras.
\end{warningbox}

\subsection{Resumen de la sección}

La fórmula de Tweedie proporciona una solución cerrada elegante para el valor medio posterior en modelos del tipo ``observación igual a media más ruido'': basta con añadir una corrección en la dirección de la función score al valor observado. En los modelos de difusión, establece un vínculo estrecho entre la función score y la operación de desenfoque (recuperación esperada de $x_0$ a partir de $x_t$), ofreciendo un puente fundamental para comprender los métodos basados en el score.
